%! Piecewise & Step Functions — Unit 1, Lesson 1.4 — Slides (Beamer)
% Deck follows the lesson's gradual-release flow (Main Ideas / Notes table shape, 2026-09-12):
% targets → warm-up → hook (left unresolved) → I do (the notes table's four instruction rows,
% one frame each, labelled by row) → we do (the Guided Practice row, problems 10–13) → debrief →
% close & start the homework. There is no group activity and no
% individual-practice launch frame: the homework IS the individual practice, started in class.
\documentclass[aspectratio=169, 11pt]{beamer}
\usepackage{algebra2-beamer}   % provides \forestheader{} and \sectionlabel[color]{}
% Coordinate grid: {xmin}{xmax}{ymin}{ymax}
\newcommand{\lgrid}[4]{%
  \draw[step=1, linegray!40, very thin] (#1,#3) grid (#2,#4);
  \draw[->, thick, charcoal] (#1,0)--(#2,0) node[below right, font=\tiny]{$x$};
  \draw[->, thick, charcoal] (0,#3)--(0,#4) node[left, font=\tiny]{$y$};}
% Months grid for the streaming service: x = months 0..7, one y-unit = $12.
\newcommand{\mgrid}{%
  \draw[step=1, linegray!40, very thin] (0,0) grid (7,4);
  \draw[->, thick, charcoal] (0,0)--(7.6,0) node[below right, font=\tiny]{$m$};
  \draw[->, thick, charcoal] (0,0)--(0,4.6) node[left, font=\tiny]{\$};
  \foreach \x in {1,...,7}{\node[below, font=\tiny, charcoal] at (\x,0){$\x$};}
  \foreach \y/\lab in {1/12,2/24,3/36,4/48}{\node[left, font=\tiny, charcoal] at (0,\y){$\lab$};}}
% Endpoint dots: {color}{x}{y}. Closed = filled; open = white-filled ring.
\newcommand{\fdot}[3]{\fill[#1] (#2,#3) circle (3.2pt);}
\newcommand{\hdot}[3]{\draw[very thick, #1, fill=white] (#2,#3) circle (3.2pt);}

\begin{document}

% TITLE SLIDE (CourseName is not defined in beamer, so the name is literal)
{
\setbeamercolor{background canvas}{bg=forest}
\begin{frame}[plain]
  \vspace{2.8cm}\hspace{0.55cm}%
  \begin{minipage}{0.9\textwidth}
    {\color{goldacc}\bfseries\small LESSON 1.4}\\[0.5cm]
    {\color{white}\bfseries\LARGE Piecewise \& Step Functions}\\[1.2cm]
    {\color{forestmid}\small Algebra 2~$\cdot$~Unit 1: Linear Functions}
  \end{minipage}
\end{frame}
}

% LEARNING TARGETS
\begin{frame}[t, plain]
  \forestheader{Today's targets}
  \sectionlabel{What you will be able to do}
  \begin{itemize}\setlength{\itemsep}{6pt}
    \item Read a \textbf{piecewise-defined function} --- one function, several rules --- and
          \textbf{evaluate} it by choosing the piece first, watching the \textbf{boundary}.
    \item Write $|x|$, and any V, as a \textbf{two-piece} linear rule.
    \item Read a piecewise graph with \textbf{open} and \textbf{closed} dots, and decide whether
          the pieces \emph{meet} or \emph{jump} (a \textbf{discontinuity}).
    \item Evaluate the \textbf{greatest-integer function} $\lfloor x\rfloor$ --- round
          \emph{down} --- and read its staircase.
  \end{itemize}
  \vspace{4pt}
  \begin{block}{How today runs}
    Warm-up (5) $\to$ notes and one function together (35) $\to$ debrief (10) $\to$
    \textbf{start the homework, on your own} (10).
  \end{block}
\end{frame}

% WARM-UP
\begin{frame}[t, plain]
  \forestheader{Warm-Up: three things you already know}
  \sectionlabel{5 minutes --- alone, then check with a neighbour}
  \begin{enumerate}\setlength{\itemsep}{7pt}
    \item \textbf{Evaluate the rule.} $g(x)=2x-1$ at $-2$, $0$, $3$; and $h(x)=5-x$ at $2$
          and $-4$.
    \item \textbf{The V is two lines.} $y=|x|$ behaves like $y=$ ? when $x\ge0$ and like $y=$ ?
          when $x<0$. Check it with $-6$.
    \item \textbf{Endpoints.} Is $1$ a solution of $x<1$? Of $x\le1$? Open or closed dot for
          $x\ge2$?
  \end{enumerate}
  \vspace{6pt}
  \begin{block}{Hold on to this}
    One rule is for $x<1$ and a different rule for $x\ge1$. The input $x=1$ fails $x<1$ and
    satisfies $x\le1$ --- so which rule gets it? Keep that question.
  \end{block}
\end{frame}

% HOOK — left unresolved
\begin{frame}[t, plain]
  \forestheader{Hook: free for three months}
  \sectionlabel[goldacc]{A streaming service, then \$12 a month}
  \begin{itemize}\setlength{\itemsep}{5pt}
    \item Total paid after 2 months? After 5? \; \textbf{\$0, \$24.}
    \item Could one straight line describe the total for every month? \; \textbf{No} --- flat,
          then rising. One function, \emph{two} rules.
  \end{itemize}
  \[ C(m)=\begin{cases} 0, & 0\le m\le 3 \\ 12(m-3), & m>3 \end{cases} \]
  \begin{block}{Vote --- and we do not settle it yet}
    \textbf{Month 3: the last free month, or the first paid one?} \quad \emph{free} \; / \;
    \emph{\$12} \quad --- Every question today is this picture: which rule does a month use, and
    what does the graph do where the rule changes?
  \end{block}
\end{frame}

% NOTES ROW 1 — SEVERAL RULES: evaluate by piece
\begin{frame}[t, plain]
  \forestheader{Notes: one function, several rules --- pick the piece first}
  \sectionlabel{notes, Several Rules $\cdot$ I do problem 1, you do problem 2}
  \begin{columns}[T]
    \begin{column}{0.52\textwidth}
      $C(m)=\begin{cases} 0, & 0\le m\le 3 \\ 12(m-3), & m>3 \end{cases}$
      \begin{itemize}\setlength{\itemsep}{4pt}
        \item $C(2)$: $2\le3$, first piece: $0$.
        \item $C(5)$: $5>3$, second piece: $12(2)=24$.
        \item $C(3)$: the \textbf{boundary}. $3$ satisfies $0\le m\le3$ --- the $\le$ includes
              it --- so $C(3)=0$. \textbf{Month 3 is still free.}
      \end{itemize}
    \end{column}
    \begin{column}{0.44\textwidth}
      \begin{block}{Boundary rule}
        An input on the boundary belongs to the piece whose inequality \emph{includes} it ---
        $\le$ or $\ge$, never the strict one. The story does not decide it; the inequality does.
      \end{block}
    \end{column}
  \end{columns}
\end{frame}

% NOTES ROW 2 — TWO PIECES
\begin{frame}[t, plain]
  \forestheader{Notes: the V is two pieces}
  \sectionlabel{notes, Two Pieces $\cdot$ Lesson 1.3's trail, written with a brace (problems 3--4)}
  \begin{columns}[T]
    \begin{column}{0.40\textwidth}
      \[ |x|=\begin{cases} x, & x\ge0 \\ -x, & x<0 \end{cases} \]
      Check: $|-3|=-(-3)=3$.
    \end{column}
    \begin{column}{0.56\textwidth}
      $H(x)=2|x-3|+1$, boundary at the vertex $x=3$:
      \begin{align*}
        x\ge3:\quad H(x) &= 2(x-3)+1 = 2x-5 \\
        x<3:\quad H(x) &= 2(3-x)+1 = 7-2x
      \end{align*}
    \end{column}
  \end{columns}
  \vspace{4pt}
  \begin{block}{Meets at the vertex}
    Right ray slope $2$, left ray slope $-2$ --- and both rules give $H(3)=1$, so the two pieces
    \emph{meet}: one connected V. That is the first ``meets''; the next row shows a ``jumps.''
  \end{block}
\end{frame}

% NOTES ROW 3 — MEETS OR JUMPS
\begin{frame}[t, plain]
  \forestheader{Notes: open or closed, meets or jumps}
  \sectionlabel{notes, Meets or Jumps $\cdot$ same story, two functions (problems 5--6)}
  \begin{columns}[T]
    \begin{column}{0.48\textwidth}
      \textbf{Total paid $C(m)$} \\[2pt]
      \centering
      \begin{tikzpicture}[scale=0.34]
        \mgrid
        \draw[very thick, forest] (0,0)--(3,0);
        \draw[very thick, forest] (3,0)--(7,4);
        \fdot{forest}{3}{0}
      \end{tikzpicture} \\[2pt]
      At $m=3$ both rules give $0$: the pieces \textbf{meet}. Constant on $[0,3]$, increasing after.
    \end{column}
    \begin{column}{0.48\textwidth}
      \textbf{Price this month $P(m)$} \\[2pt]
      \centering
      \begin{tikzpicture}[scale=0.34]
        \mgrid
        \draw[very thick, redacc] (0,0)--(3,0);
        \draw[very thick, redacc] (3,1)--(7,1);
        \fdot{redacc}{3}{0}
        \hdot{redacc}{3}{1}
      \end{tikzpicture} \\[2pt]
      At $m=3$ the rules give $0$ and $12$: the graph \textbf{jumps}. Range $\{0,12\}$.
    \end{column}
  \end{columns}
  \vspace{3pt}
  \begin{block}{One test}
    Both rules agree at the boundary, or they do not. Same value: meet. Different: jump --- a
    \textbf{discontinuity}. Closed dot: the endpoint is \emph{in}; open dot: it is \emph{out}.
  \end{block}
\end{frame}

% NOTES ROW 4 — BOUNDARY: the crux (problem 7), then rounding down (problem 9)
\begin{frame}[t, plain]
  \forestheader{Notes: exactly one piece owns the boundary}
  \sectionlabel[redacc]{notes, Boundary $\cdot$ the crux, problem 7 --- the story cannot decide it; the inequality does}
  \begin{itemize}\setlength{\itemsep}{5pt}
    \item ``$P(3)=12$, because the price is \$12 after the free months''? \textbf{No.} $3$
          satisfies $0\le m\le3$, so $P(3)=0$ --- the \emph{closed} dot at $(3,0)$ says so.
    \item The open dot at $(3,12)$ is \emph{not} a point of the graph. It marks where the second
          piece starts \emph{without} its endpoint. Read the wrong dot and you have given the
          boundary to the wrong piece.
  \end{itemize}
  \vspace{3pt}
  \begin{block}{Step functions: rounding down}
    $P(m)$ is constant, then jumps --- a \textbf{step function}. $\lfloor x\rfloor$ = the greatest
    integer $\le x$: $\lfloor3.7\rfloor=3$, $\lfloor5\rfloor=5$, $\lfloor-0.4\rfloor=-1$,
    $\lfloor-2.5\rfloor=-3$ --- down means \emph{more negative}. The staircase is constant $n$ on
    $[n,n+1)$: closed left, open right --- the boundary rule, one step at a time. (Problem 9.)
  \end{block}
\end{frame}

% GUIDED PRACTICE — WE DO
\begin{frame}[t, plain]
  \forestheader{Guided Practice: we do this one together}
  \sectionlabel[goldacc]{You hold the pen --- problems 10--13: one function, then change one piece}
  $p(x)=-2x$ for $x\le-1$, \; $x+3$ for $x>-1$ --- the graph is on your page.
  \begin{enumerate}\setlength{\itemsep}{4pt}
    \item[\textbf{10.}] Evaluate at $-3$, $0$, $2$ --- then at $-1$: which piece owns it, and
          how do you know?
    \item[\textbf{11.}] Meets or jumps at $-1$? Decreasing where, increasing where, and the
          minimum?
    \item[\textbf{12.}] Change the second piece to $x+4$: who owns $-1$ now, meet or jump, which
          dot is closed? A classmate said ``both rules gave $2$, so it does not matter who owns
          it'' --- what did they miss?
    \item[\textbf{13.}] Round down: $\lfloor2.9\rfloor$, $\lfloor-1.2\rfloor$, $\lfloor4\rfloor$
          --- and does the step $[4,5)$ own $5$?
  \end{enumerate}
  \vspace{4pt}
  \begin{block}{The questions I will ask}
    ``Does $-1$ satisfy $x\le-1$? Does it satisfy $x>-1$?'' \quad ``You changed the rule --- did
    you change the \emph{inequality}?'' \quad ``Which integer is just \emph{below} $-1.2$?''
  \end{block}
\end{frame}

% DEBRIEF
\begin{frame}[t, plain]
  \forestheader{Debrief: pens down, papers up}
  \sectionlabel[redacc]{10 minutes --- aloud, nothing to fill in}
  \begin{enumerate}\setlength{\itemsep}{5pt}
    \item \textbf{$C$ and $P$ side by side.} Same value at $m=3$: meet. Different: jump. Point at
          the two dots that earned it.
    \item \textbf{Problem 12.} The first piece still owns $-1$ --- $p(-1)=2$, closed dot at
          $(-1,2)$, open dot at $(-1,3)$, and the graph jumps. Whoever moved $p(-1)$ to $3$ read
          the open dot. Same as $P(3)=0$.
    \item \textbf{Rounding down.} $\lfloor-1.2\rfloor=-2$, not $-1$. And $\lfloor5\rfloor=5$:
          the step $[4,5)$ stops \emph{without} $5$; the next one starts \emph{with} it.
    \item \textbf{Pieces can meet.} The trail $H(x)$ and $p(x)$ both do --- ``piecewise'' does not
          mean ``discontinuous.''
  \end{enumerate}
  \vspace{4pt}
  \begin{block}{Say it without the notes}
    Month 3: $3$ satisfies $0\le m\le3$, so the first piece owns it, so $P(3)=0$ --- closed dot at
    $(3,0)$.
  \end{block}
\end{frame}

% CLOSE & START HOMEWORK
\begin{frame}[t, plain]
  \forestheader{Close --- and start the homework, on your own}
  \sectionlabel{10 minutes $\cdot$ silent $\cdot$ the Homework page in your packet}
  You walked in with functions that follow \emph{one} rule. You walk out able to read a function
  that follows \emph{several} --- pick the piece, mind the boundary, read the dots --- and to see
  Lesson 1.3's V as the first piecewise function you already knew.
  \vspace{5pt}
  \begin{block}{Homework --- scored, due the first class after two study halls}
    We work item~1's first blank together right now; then you keep going, alone. Six items:
    evaluating with a boundary, a meets-or-jumps pair, a graph to read \emph{and} write as a rule,
    the staircase, a phone plan, and one multiple-choice item. Whatever is left is due the first
    class after two study halls --- I will announce the date in class and on TurtleNet.
  \end{block}
  \vspace{4pt}
  \textbf{Next:} Lesson 1.5 --- we leave exact rules behind and fit one \emph{line of best fit} to
  scattered data, then measure how well it fits.
\end{frame}

\end{document}
